% 2. Pregled oblasti i srodni radovi
\section{Pregled oblasti i srodni radovi}
\label{sec:pregled}

Primena veštačke inteligencije u obrazovanju razvijala se od rane automatizacije zasnovane na pravilima ka adaptivnim sistemima koji tumače i generišu složene odgovore nalik ljudskim.
Pod veštačkom inteligencijom se podrazumevaju računarski sistemi projektovani da obavljaju kognitivne funkcije poput učenja, rezonovanja i rešavanja problema~\cite{baker2019educ},
a u obrazovnom kontekstu skup tehnologija koje omogućavaju personalizovanija i adaptivnija okruženja za učenje~\cite{zawacki2019systematic}.
Sistematski pregledi im pripisuju potencijal da unaprede ishode učenja, prošire dostupnost obrazovanja i efikasnije iskoriste nastavne resurse~\cite{yuskovych2022application,zhang2023study}.
Veliki jezički modeli tu donose promenu u odnosu na ranije pristupe zasnovane na obradi prirodnog jezika:
ne zahtevaju ručno projektovane atribute ni obučavanje za konkretan domen,
već obrađuju slobodan tekst i rezonuju o značenju na višem nivou apstrakcije.
Zato se ispituju za nastavu usmerenu na učenika, automatsko generisanje sadržaja i nove oblike provere znanja~\cite{zhang2023study}.

Empirijske studije ocenjivanja uz podršku velikih jezičkih modela daju mešovite rezultate.
Kuli i Jusuf~\cite{kooli2025transforming} ispitali su ocenjivanje uz pomoć ChatGPT-a na 25 studentskih radova iz kohorte od približno 100 studenata druge godine na kursu iz društvenih nauka,
gde su radove nezavisno ocenili i ljudski ocenjivači i model.
Dobijena je umerena korelacija između ocena veštačke inteligencije i ljudskih ocena (koeficijent korelacije približno 0{,}45): model aproksimira opšte obrasce ocenjivanja,
ali pojedinačne evaluacije se i dalje razlikuju.
Autori ističu da ChatGPT ubrzava ocenjivanje i daje korisnu povratnu informaciju,
ali da mu izostaje kontekstualno razumevanje nijansiranih studentskih argumenata.
Upozoravaju i na rizik produbljivanja obrazovnih nejednakosti i na potrebu za institucionalnim okvirima za transparentnost, nadzor i ublažavanje pristrasnosti pri uvođenju takvih alata.

Lundgren~\cite{lundgren2024large} je ispitao GPT-4 na 60 eseja master nivoa iz političkih nauka koje su prethodno ocenili univerzitetski nastavnici.
Prosečne ocene modela bile su usklađene sa nastavničkim, ali je pouzdanost među ocenjivačima (između modela i ljudi) bila niska:
model prati agregirane trendove, a ne fino ljudsko rasuđivanje.
Ocenjivanje je bilo sklono izbegavanju rizika i oslanjanju na opšte karakteristike eseja, poput kvaliteta jezika, umesto na detaljne rubrike.
Izmene prompta i uputstava za ocenjivanje donele su samo ograničena poboljšanja,
pa su tadašnji sistemi bili slabo osetljivi na suptilne varijacije u kriterijumima evaluacije.

Jukjevič~\cite{jukiewicz2024future} je ChatGPT primenio na ocenjivanje devet programerskih zadataka na petnaestonedeljnom kursu programiranja u Python-u sa 67 studenata osnovnih studija kognitivnih nauka,
gde su radove nezavisno ocenili i nastavnik i model.
Korelacija sa nastavničkim ocenama bila je snažno pozitivna, mada je model dodeljivao nešto strože ocene.
Ponovljivost kroz višestruka pokretanja bila je visoka, sa zanemarljivim varijacijama između iteracija.
Uz numeričku ocenu, ChatGPT je procenjivao ispravnost koda, proveravao poštovanje smernica PEP8 i za nekoliko sekundi generisao detaljnu povratnu informaciju.
Autor navodi i ograničenja: halucinirane greške, troškove korišćenja API-ja i povremene nedoslednosti u ocenjivanju koje zahtevaju proveru nastavnika.
Zaključak studije je da je ChatGPT koristan komplementaran alat za procenu koda i generisanje povratnih informacija,
ali da ljudski nadzor ostaje neophodan za pouzdanost i pravednost.

Ovi nalazi se uklapaju u širu sliku: pouzdanost ocenjivanja velikim jezičkim modelima zavisi od tipa zadatka.
Slaganje sa ljudskim ocenjivačima veće je tamo gde su odgovori strukturirani, kao kod definicija, formula ili fragmenata koda,
a opada kod subjektivnih ili otvorenih pisanih zadataka~\cite{tomic2022ai,vittorini2021ai}.
U radu~\cite{pinto2023large} GPT-3 je evaluiran na odgovorima otvorenog tipa profesionalaca iz industrije i pokazao snažne korelacije sa ekspertskim ocenama, pod uslovom da su kriterijumi jasno definisani i kontekstualno utemeljeni.
Izvan numeričkog ocenjivanja, sistemi zasnovani na veštačkoj inteligenciji generišu adaptivne aktivnosti učenja,
isporučuju personalizovanu povratnu informaciju i podržavaju formativno i sumativno ocenjivanje prilagođeno individualnim karakteristikama učenika~\cite{diwan2023ai,luckin2016intelligence}.
Inteligentni tutorski sistemi i srodni pristupi nastoje da aproksimiraju individualnu nastavnu podršku uz veću doslednost i skalabilnost evaluacije~\cite{conati2009intelligent,garcia2018social}.
Kritičke perspektive, s druge strane, upozoravaju na „zamagljivanje'' odgovornosti u automatskom ocenjivanju,
gde odluke postaju neprozirne, a odgovornost razvodnjena~\cite{swiecki2022assessment}.
Transparentnost, objašnjivost, privatnost podataka i kvalitet povratnih informacija zato ostaju centralne teme konceptualnih analiza ocenjivanja podržanog veštačkom inteligencijom~\cite{kumar2023faculty}.

Kroz ovu literaturu provlači se isti zaključak: veliki jezički modeli su pomoćni alati koji automatizuju rutinske aspekte ocenjivanja,
dok nastavnik ostaje nezamenljiv za pedagoško rasuđivanje, kontekstualno tumačenje i granične slučajeve~\cite{van2022applications,zawacki2019systematic}.
Preporukama iz literature najbliže odgovaraju pristupi koji spajaju automatizaciju sa nadzorom čoveka u petlji, podesivom strogošću ocenjivanja i interpretabilnom povratnom informacijom.
Ta tri zahteva, kontrola nastavnika, transparentnost i prilagodljivost, oblikovala su i platformu EduGrader predstavljenu u ovom radu.
